% ===================== PREAMBULE COMMUN (à copier VERBATIM en tête de CHAQUE .tex) =====================
\documentclass[11pt,a4paper]{article}
\usepackage[utf8]{inputenc}
\usepackage[T1]{fontenc}
\usepackage{lmodern}
\usepackage[french]{babel}
\usepackage{amsmath,amssymb,mathtools}
\usepackage{icomma}
\usepackage{xcolor}
\usepackage{colortbl}
\usepackage{tikz}
\usepackage{pgfplots}
\usepackage[most]{tcolorbox}
\usepackage{enumitem}
\usepackage{array}
\usepackage{booktabs}
\usepackage{multicol}
\usepackage{fancyhdr}
\usepackage{caption}
\captionsetup{labelformat=empty,justification=centering,font=small}
\usepackage[hidelinks]{hyperref}
\usetikzlibrary{arrows.meta,positioning,calc,angles,quotes,patterns,backgrounds,shapes.geometric}
\pgfplotsset{compat=1.18}
\usepackage[a4paper,margin=2.2cm,headheight=26pt,headsep=10pt]{geometry}

\definecolor{primary}{RGB}{219,54,77}
\definecolor{primarydark}{RGB}{150,28,46}
\definecolor{primarylight}{RGB}{250,224,228}
\definecolor{defcol}{RGB}{0,114,178}
\definecolor{thmcol}{RGB}{0,140,120}
\definecolor{excol}{RGB}{0,135,90}
\definecolor{methcol}{RGB}{90,90,100}
\definecolor{courbe}{RGB}{40,90,200}
\definecolor{courbeB}{RGB}{210,60,40}
\definecolor{courbeC}{RGB}{30,150,80}
\definecolor{axiscol}{RGB}{60,60,60}
\definecolor{gridcol}{RGB}{200,205,215}

\tcbset{boxbase/.style={enhanced,breakable,boxrule=0.9pt,arc=2.5pt,
   left=3mm,right=3mm,top=2.3mm,bottom=2.3mm,
   fonttitle=\bfseries,coltitle=white,toptitle=0.6mm,bottomtitle=0.6mm}}
\newtcolorbox{methode}[1][]{boxbase,colback=methcol!7,colframe=methcol,sharp corners,
  title={\textsc{Méthode}\ifx\relax#1\relax\relax\else\textmd{ --- \itshape #1}\fi}}
\newtcolorbox{exemple}[1][]{boxbase,colback=excol!5,colframe=excol,
  title={\textsc{Exemple}\ifx\relax#1\relax\relax\else\textmd{ : \itshape #1}\fi}}
\newtcolorbox{idee}[1][]{boxbase,colback=defcol!6,colframe=defcol,
  title={\textsc{L'essentiel}\ifx\relax#1\relax\relax\else\textmd{ : \itshape #1}\fi}}
\newtcolorbox{piege}[1][]{boxbase,colback=primary!7,colframe=primary,
  title={\textsc{Piège}\ifx\relax#1\relax\relax\else\textmd{ : \itshape #1}\fi}}

\pgfplotsset{repere/.style={axis lines=middle,
    axis line style={-{Stealth[length=2.2mm]},axiscol,line width=0.8pt},
    label style={font=\footnotesize},tick label style={font=\scriptsize},
    every axis x label/.style={at={(current axis.right of origin)},anchor=north west},
    every axis y label/.style={at={(current axis.above origin)},anchor=south east},
    xlabel={$x$},ylabel={$y$},grid=both,minor tick num=0,
    major grid style={gridcol,line width=0.4pt},samples=200,clip=false}}

\pagestyle{fancy}\fancyhf{}
\fancyhead[L]{\small\textcolor{primarydark}{\textbf{Fiche 4} — Les fonctions}}
\fancyhead[R]{\small\textcolor{primarydark}{\textsc{Carnet d'entraînement}}}
\fancyfoot[C]{\small\thepage}
\renewcommand{\headrulewidth}{0.8pt}
\renewcommand{\headrule}{\hbox to\headwidth{\color{primary}\leaders\hrule height \headrulewidth\hfill}}
\usepackage{titlesec}
\titleformat{\section}{\Large\bfseries\color{primarydark}}{\thesection}{0.6em}{}
\titleformat{\subsection}{\large\bfseries\color{primary}}{\thesubsection}{0.6em}{}

\newcommand{\R}{\mathbb{R}}\newcommand{\Z}{\mathbb{Z}}\newcommand{\N}{\mathbb{N}}
\newcommand{\Df}{D_f}
\newcommand{\croit}{\textcolor{thmcol}{\Large$\nearrow$}}
\newcommand{\decroit}{\textcolor{thmcol}{\Large$\searrow$}}

\newcommand{\exercice}[1]{\medskip\noindent{\color{primarydark}\bfseries\large Exercice #1}\par\nopagebreak\smallskip}
\newcommand{\titrecarnet}[3]{%
\begin{center}\begin{tikzpicture}
\node[rectangle,rounded corners=7pt,fill=primary,text=white,minimum width=15cm,
  minimum height=1.7cm,align=center,inner sep=8pt]
  {{\LARGE\bfseries #1}\\[2pt]{\normalsize #2}};
\end{tikzpicture}\end{center}
\begin{center}\itshape\color{primarydark}#3\end{center}\medskip}

% ---- RÈGLES D'USAGE DES MACROS ----
% * \croit et \decroit s'emploient en mode TEXTE (dans une cellule de tableau),
%   JAMAIS entre $...$ (ils contiennent déjà un $\nearrow$).
% * \titrecarnet{Titre}{Sous-titre}{Ligne en italique} : bandeau de titre.
% * \exercice{1} : titre d'exercice.
% * Boîtes : \begin{idee}...\end{idee}, \begin{exemple}[titre]...\end{exemple},
%   \begin{methode}[titre]...\end{methode}, \begin{piege}[titre]...\end{piege}.
% Le corps du document commence après \begin{document}.
\begin{document}
\titrecarnet{Thème 3 — Géométrie analytique des droites}{Corrigé — niveau Difficile}{Détail des calculs.}

\exercice{1}
\begin{enumerate}[label=\alph*),leftmargin=1.8em,topsep=2pt]
  \item $d$ a pour pente $2$ ; $d'\perp d$ a pour pente $a'=-\dfrac12$. Comme $d'$ passe par $(0;4)$, son ordonnée à l'origine est $4$ :
        \[ \boxed{d':\ y=-\tfrac12 x+4}. \]
  \item $y=0\Rightarrow -\dfrac12 x+4=0\Rightarrow x=8$. Abscisse cherchée : $\boxed{8}$.
\end{enumerate}

\exercice{2}
\begin{enumerate}[label=\alph*),leftmargin=1.8em,topsep=2pt]
  \item $\overline{AB}=\sqrt{(4-0)^2+(0-0)^2}=4$ ; \quad $\overline{BC}=\sqrt{(4-4)^2+(3-0)^2}=3$ ; \quad $\overline{AC}=\sqrt{(4-0)^2+(3-0)^2}=\sqrt{25}=5$.
  \item $\overline{AB}^2+\overline{BC}^2=16+9=25=\overline{AC}^2$. La réciproque de Pythagore donne un angle droit en $B$.
  \item Les côtés de l'angle droit sont $\overline{AB}=4$ et $\overline{BC}=3$, donc aire $=\dfrac{4\times 3}{2}=\boxed{6}$.
\end{enumerate}

\exercice{3}
\begin{enumerate}[label=\alph*),leftmargin=1.8em,topsep=2pt]
  \item Pente de $d_1=(AB)$ : $\dfrac{6-3}{7-2}=\dfrac35$. Perpendiculaire : $a_2=-\dfrac{5}{3}$. $d_2$ passe par $C(4;1)$ :
        $b=1-\big(-\tfrac53\big)\times 4=1+\dfrac{20}{3}=\dfrac{23}{3}$. Donc $\boxed{d_2:\ y=-\tfrac53 x+\tfrac{23}{3}}$.
  \item On résout $-\dfrac53 x+\dfrac{23}{3}=-2x+8$. En multipliant par $3$ : $-5x+23=-6x+24\Rightarrow x=1$. Puis $y=-2\times 1+8=6$.
        \[ \boxed{(1;6)}. \]
\end{enumerate}

\exercice{4}
\begin{enumerate}[label=\alph*),leftmargin=1.8em,topsep=2pt]
  \item Branche droite $(AC)$ avec $A\big(\tfrac32;0\big)$, $C(3;3)$ : $a=\dfrac{3-0}{3-\tfrac32}=\dfrac{3}{\tfrac32}=2$ ; $b=3-2\times 3=-3$. Donc $\boxed{y=2x-3}$.
  \item Branche gauche $(AB)$ avec $A\big(\tfrac32;0\big)$, $B(0;3)$ : $a=\dfrac{3-0}{0-\tfrac32}=-2$ ; $b=3$. Donc $\boxed{y=-2x+3}$.
        (On retrouve bien $f(x)=\lvert 2x-3\rvert$.)
  \item $\overline{AB}=\sqrt{\big(0-\tfrac32\big)^2+(3-0)^2}=\sqrt{\tfrac94+9}=\sqrt{\tfrac{45}{4}}=\dfrac{3\sqrt5}{2}$.\\
        $\overline{AC}=\sqrt{\big(3-\tfrac32\big)^2+(3-0)^2}=\sqrt{\tfrac94+9}=\dfrac{3\sqrt5}{2}$.\\
        Comme $\overline{AB}=\overline{AC}$, le triangle $ABC$ est \textbf{isocèle en $A$}. \checkmark
  \item Forme canonique de $(AB)$ : $y=-2x+3\Rightarrow 2x+y-3=0$, donc $a=2,\ b=1,\ c=-3$.
        \[ \operatorname{dist}(P,AB)=\frac{\lvert 2\times(-2)+1\times 1-3\rvert}{\sqrt{2^2+1^2}}=\frac{\lvert -6\rvert}{\sqrt5}=\boxed{\dfrac{6}{\sqrt5}}. \]
\end{enumerate}
\end{document}
